% Part: applied-modal-logics
% Chapter: epistemic-logic
% Section: introduction

\documentclass[../../../include/open-logic-section]{subfiles}

\begin{document}

\olfileid{aml}{el}{int}

\olsection{مقدمه}

همان‌گونه که منطق موجهات به \emph{گزاره‌های موجهاتی} و روابط استلزام
میان آن‌ها می‌پردازد، منطق معرفتی به \emph{گزاره‌های معرفتی} و روابط
استلزام میان آن‌ها می‌پردازد. به‌جای تفسیر عملگرهای موجهاتی به‌عنوان
نمایانگر امکان و ضرورت، رابط‌های یک‌موضعی به شیوه‌های معرفتی یا باوری
تفسیر می‌شوند تا دانش و باور را مدل‌سازی کنند. برای نمونه، شاید بخواهیم
ادعاهایی مانند موارد زیر را بیان کنیم:

\begin{enumerate}
\item ریچارد می‌داند کلگری در آلبرتاست.
\item آدری گمان می‌کند ممکن است سگی روی مبل باشد.
\item ریچارد می‌داند آدری می‌داند کلاسش سه‌شنبه‌ها برگزار می‌شود.
\item همه می‌دانند یک سال ۱۲ ماه دارد.
\end{enumerate}
منطق معرفتی معاصر را غالباً به کتاب \emph{دانش و باور} یاکو هینتیکا
(Jaakko Hintikka) در سال ۱۹۶۲ بازمی‌گردانند؛ کتابی که در زمانی نوشته
شد که معناشناسی جهان‌های ممکن در منطق کاربرد روزافزونی می‌یافت. در واقع،
منطق‌های معرفتی بیشتر ابزارهای معناشناختیِ منطق‌های موجهاتی دیگر را به
کار می‌برند، اما آن‌ها را متفاوت تفسیر می‌کنند. تغییر اصلی در چیزی است
که \emph{رابطهٔ دسترسی‌پذیری} را نمایانگرش می‌گیریم. در منطق‌های معرفتی،
این روابط گونه‌ای \emph{امکان معرفتی} را بازنمایی می‌کنند. خواهیم دید
مفهوم معرفتی‌ای که مدل می‌کنیم بر قیودی اثر می‌گذارد که می‌خواهیم بر
رابطهٔ دسترسی‌پذیری بگذاریم. همچنین خواهیم دید وقتی نظریهٔ تناظر تفسیری
معرفتی می‌یابد چه رخ می‌دهد. در نمونه‌های بالا دو عامل، ریچارد و آدری،
و رابطهٔ میان چیزهایی که هریک می‌داند ذکر شده‌اند. منطق‌های معرفتی مورد
بررسی ما چندعاملی خواهند بود تا بتوان چنین اموری را در آن‌ها بیان کرد.
در مقابل، منطق معرفتی تک‌عاملی تنها دربارهٔ آنچه یک فرد می‌داند یا
باور دارد سخن می‌گوید.

\end{document}
